% NeurIPS 2026 paper checklist.
%
% Kept in a SEPARATE file on purpose. The camera-ready workshop limit is seven
% content pages
% against the body only; the checklist is submitted alongside and must not be
% allowed to push the body over. Compile standalone with build_checklist.sh, or
% \input it after \bibliography when the venue asks for a single file.
%
% Every answer below points at the section that actually contains the evidence.
% "Yes" is used only where the PAPER reports the thing, not where we merely
% could have. Where evidence exists only in artifacts or code and not in the
% manuscript, that is stated explicitly in the justification.

\begin{enumerate}

\item {\bf Claims}
    \item[] Question: Do the main claims made in the abstract and introduction
    accurately reflect the paper's contributions and scope?
    \item[] Answer: \answerYes{}
    \item[] Justification: The abstract and the three contribution bullets
    (\S\ref{sec:intro}) each name the measurement that backs them: the
    granularity sweep in three adapter-construction regimes
    (\S\ref{sec:curve}), the storage/construction/routing accounting
    (\S\ref{sec:cost}), and the five-way operating-point comparison including
    full context and top-5 retrieval (\S\ref{sec:operating}). Scope is stated in
    the same sentences: one 4B model, one \NPerArm{}-fact TOFU memory,
    per-adapter rank fixed at 16. Where the abstract says the monotonic
    original-question decline is not observed under equal per-fact exposure, the
    text states that this is an unresolved comparison rather than an equivalence
    result (\S\ref{sec:curve}). The abstract closes by declining to name a
    winner, which matches \S\ref{sec:operating}.

\item {\bf Limitations}
    \item[] Question: Does the paper discuss the limitations of the work
    performed by the authors?
    \item[] Answer: \answerYes{}
    \item[] Justification: \S\ref{sec:limits} is four dedicated paragraphs:
    routing is unusually easy on TOFU and our routing numbers are an upper
    bound; a single base model, target-module set, rank and learning rate, with
    no capacity sweep, so no claim about how these curves move with rank; the
    construction budget is part of the result, so a granularity sweep does not
    measure adapter capacity on its own and we tested two points in the
    budget/exposure space rather than the space itself, with scheduler horizon
    not separated from exposure in the equal-exposure arm; and the in-context
    comparison charges top-5 retrieval for prompt tokens but not for index,
    corpus, or retriever infrastructure, does not vary corpus scale (citing a
    paper that reports the opposite ordering at a smaller corpus), used one
    fixed full-context ordering so position is confounded with item difficulty,
    and measures factual recall only, since TOFU holds authorial voice constant.
    The routing tax is explicitly reported as unresolved rather than zero,
    because it falls below the engine's own reproducibility floor
    (\S\ref{sec:cost}).

\item {\bf Theory assumptions and proofs}
    \item[] Question: For each theoretical result, does the paper provide the
    full set of assumptions and a complete (and correct) proof?
    \item[] Answer: \answerNA{}
    \item[] Justification: The paper is purely empirical. It states no theorem
    and derives no analytical result.

\item {\bf Experimental result reproducibility}
    \item[] Question: Does the paper fully disclose all the information needed
    to reproduce the main experimental results?
    \item[] Answer: \answerYes{}
    \item[] Justification: \S\ref{sec:design} gives the dataset and exact config
    (TOFU \texttt{forget10\_perturbed}), the partition rule (author-aligned
    20-fact blocks packed in ascending order, which is deterministic), the base
    model, rank, learning rate, $\alpha$, target module, optimizer, batch size
    and gradient accumulation, the grading rule and threshold, the pairing and
    test procedure, and the serving stack with versions. Both gradient
    construction regimes are specified separately: the fixed-update regime
    (exactly \SeedSteps{} optimizer updates per adapter, RNGs reset to seed
    \SeedSeed{} immediately before LoRA initialization, termination verified by
    reading \texttt{global\_step} back) and the equal-exposure regime
    (\ExpEpochs{} epochs per partition, trained independently each with its own
    scheduler horizon, giving \ExpGTwentySteps{}--\ExpGTwoHundredSteps{}
    updates). Two controls needed to reproduce the result honestly are also
    described: the per-document truncation gate (\S\ref{sec:design}) and the
    log-probability activity gate (\S\ref{sec:reliability}). Caveats stated in
    the paper: the serving stack is not bitwise reproducible under concurrency
    (\S\ref{sec:cost}), so exact per-item replication is not expected, and the
    full-context arm required a server configured with a 16{,}384-token context
    window rather than the configuration used by the other arms.

\item {\bf Open access to data and code}
    \item[] Question: Does the paper provide open access to the data and code?
    \item[] Answer: \answerYes{}
    \item[] Justification: The camera-ready release includes the experiment
    scripts, frozen per-item results, generated tables and figures, environment
    details, and step-by-step reproduction instructions. The public TOFU
    dataset and Qwen base model are cited at their points of use.

\item {\bf Experimental setting/details}
    \item[] Question: Does the paper specify all the training and test details
    needed to understand the results?
    \item[] Answer: \answerYes{}
    \item[] Justification: \S\ref{sec:design} specifies rank 16, the MLP-side
    target module (\texttt{down\_proj}), the learning rate ($\LearningRate{}$)
    and $\alpha$ (\LoraAlpha{}) held constant across the ladder, the oracle and
    predicted routing procedures, the grading threshold of 0.7, and
    $n=\NGradable{}$ of \NPerArm{} after dropping ungradable items pairwise. The
    two granularity regimes are distinguished explicitly, since optimizer
    updates and per-fact exposure cannot both be held fixed while partition size
    varies: fixed-update (\SeedSteps{} updates per adapter) and equal-exposure
    (\ExpEpochs{} epochs per partition). The operating-point comparison of
    \S\ref{sec:operating} uses a separate step-matched gradient arm
    (\SgdOptSteps{} steps per adapter), and the paper states that its numbers
    are not interchangeable with the granularity regimes.

\item {\bf Experiment statistical significance}
    \item[] Question: Does the paper report error bars suitably and correctly
    defined or other appropriate information about the statistical significance
    of the experiments?
    \item[] Answer: \answerYes{}
    \item[] Justification: Because the same \NPerArm{} items recur at every
    point on the ladder, comparisons are paired at the item level and tested
    with exact McNemar on discordant pairs. The end-to-end $20\to\NPerArm{}$
    paired $p$-value is reported for each of the three construction regimes, on
    both original and paraphrased questions, in the table inside
    Figure~\ref{fig:granularity}, and lost/gained counts are given in the text
    (\S\ref{sec:design}, \S\ref{sec:curve}). Interval types are distinguished:
    granularity-curve aggregate intervals use an item-clustered bootstrap that
    resamples items and carries all of an item's measurements together, whereas
    single-arm binomial intervals, including the full-context 95\% interval in
    \S\ref{sec:operating}, are Wilson intervals.
    Figure~\ref{fig:granularity} plots point estimates
    without error bars, with significance carried by the paired tests in the
    accompanying table rather than by interval overlap. The paper also states
    where a difference is \emph{not} resolvable rather than reporting it as an
    effect: the 100-to-\NPerArm{} compiled step
    ($p=\StepHundredtoTwoHundredP{}$), every equal-exposure original-question
    comparison against 20 facts per adapter (all $p\ge\ExpQPMin{}$, described as
    unresolved rather than as equivalence), and the routing tax, which sits
    below the engine's reproducibility floor.

\item {\bf Experiments compute resources}
    \item[] Question: Does the paper provide sufficient information on the
    computer resources needed to reproduce the experiments?
    \item[] Answer: \answerYes{}
    \item[] Justification: \S\ref{sec:design} states one MI300X GPU running vLLM
    0.24.0 on ROCm 7.2.3 with a 4B base model. Construction cost, which is the
    dominant compute term and is itself a reported result, is given per adapter
    in the prose of \S\ref{sec:operating} (\DtwoLBuildPer{}\,s compiled against
    \SgdBuildPer{}\,s gradient-trained), whereas Table~\ref{tab:operating}
    reports the corresponding ten-adapter bank totals (\GTwentyBuild{}\,s
    compiled against \SgdBuildTotal{}\,s gradient-trained); per-granularity
    compiled build seconds appear in the table
    inside Figure~\ref{fig:granularity}. Evaluation cost is not separately
    tabulated, and we report prompt tokens rather than wall-clock time for the
    in-context arms, since the full-context arm ran under a different server
    configuration and its runtime is not comparable to the others.

\item {\bf Code of ethics}
    \item[] Question: Does the research conducted in the paper conform, in every
    respect, with the NeurIPS Code of Ethics?
    \item[] Answer: \answerYes{}
    \item[] Justification: The study uses a public synthetic benchmark of
    fictitious author biographies and a public base model. No human subjects, no
    personal or private data, and no annotator labour are involved.

\item {\bf Broader impacts}
    \item[] Question: Does the paper discuss both potential positive societal
    impacts and negative societal impacts of the work performed?
    \item[] Answer: \answerNA{}
    \item[] Justification: The work is a measurement study of where a personal
    agent should store facts. It introduces no new model, capability, or
    artifact for release. We note that the setting concerns private user data,
    but our results are about how finely a memory is partitioned and how the
    adapters are constructed, not about whether stored information can be
    recovered or removed; the paper makes no privacy or unlearning claim, and we
    do not consider the measurement itself to carry a distinct societal risk.

\item {\bf Safeguards}
    \item[] Question: Does the paper describe safeguards that have been put in
    place for responsible release of data or models that have a high risk for
    misuse?
    \item[] Answer: \answerNA{}
    \item[] Justification: No data, model, or scraped corpus is released with
    this submission, and the adapters trained here encode only fictitious
    biographies from a public benchmark, so there is no misuse surface to guard.

\item {\bf Licenses for existing assets}
    \item[] Question: Are the creators or original owners of assets used in the
    paper properly credited and are the license and terms of use explicitly
    mentioned and properly respected?
    \item[] Answer: \answerYes{}
    \item[] Justification: Every asset used is cited at the point of use in
    \S\ref{sec:design}: TOFU~\citep{tofu} for the data, Qwen3-4B-Instruct-2507
    base~\citep{qwen3,qwen3-4b-2507} for the model, the document-to-LoRA
    hypernetwork~\citep{d2l} for the compiled generator, LoRA~\citep{lora} for
    the gradient arm, and retrieval-augmented generation~\citep{rag} for the
    retrieval arm. Note honestly: individual license names and versions are not
    printed in the manuscript; all assets used are publicly released and used
    within their published terms, and license details will accompany the code
    release.

\item {\bf New assets}
    \item[] Question: Are new assets introduced in the paper well documented and
    is the documentation provided alongside the assets?
    \item[] Answer: \answerNA{}
    \item[] Justification: This submission releases no new asset. The adapters
    built for the study are experimental artifacts, not a contributed model.

\item {\bf Crowdsourcing and research with human subjects}
    \item[] Question: For crowdsourcing experiments and research with human
    subjects, does the paper include the full text of instructions given to
    participants and screenshots, if applicable, as well as details about
    compensation?
    \item[] Answer: \answerNA{}
    \item[] Justification: No crowdsourcing and no human subjects. The
    paraphrases used in the evaluation are the human-written paraphrases shipped
    with the TOFU benchmark, not collected by us.

\item {\bf Institutional review board (IRB) approvals or equivalent for research
with human subjects}
    \item[] Question: Does the paper describe potential risks incurred by study
    participants, whether such risks were disclosed to the subjects, and whether
    Institutional Review Board (IRB) approvals (or an equivalent approval/review
    based on the requirements of your country or institution) were obtained?
    \item[] Answer: \answerNA{}
    \item[] Justification: No study participants and no human-subjects research,
    so no review was required.

\item {\bf Declaration of LLM usage}
    \item[] Question: Does the paper describe the usage of LLMs if it is an
    important, original, or non-standard component of the core methods in this
    research?
    \item[] Answer: \answerYes{}
    \item[] Justification: Language models are the object of study, and every
    such use is described in \S\ref{sec:design}: a Qwen3-4B-Instruct-2507 base as the
    model under test, a document-to-LoRA hypernetwork as the compiled-memory
    generator, gradient-trained LoRA for the gradient arm, and the same base
    prompted with
    retrieved or full-corpus context for the in-context arms. Grading is a
    deterministic token-recall rule against the gold answer, not an LLM judge.

\end{enumerate}
